% !PRINTED 451
\weberpage{451}
%<<CONT
variables (with the inclusion of the constants), in which one may restrict the constant coefficients to an arbitrary number field, say to that of the rational numbers, or may also take them entirely unrestricted.

We do not wish to multiply examples further here, since the exact investigation of fields of special kinds will be one of the principal tasks of our later expositions.

\begin{center}
\S.\,140.

\gesperrt{Adjunction.}
\end{center}

If to a field of any magnitudes whatever, which we shall denote by $\Omega$, there is added any magnitude $\alpha$ that is not contained in it, there arises a new system of magnitudes
\[
\Omega,\ \alpha,
\]
which, however, will not be a field; in order to derive a field from it, one must add all magnitudes that can be derived by combining $\alpha$ with the numbers of $\Omega$ by means of the fundamental operations of arithmetic. One thus obtains an extended field $\Omega'$, which contains both $\alpha$ and $\Omega$, and which is completely determined by $\alpha$ and $\Omega$. We shall call it the field $\Omega, \alpha$.

This adding of a new magnitude $\alpha$ to a field $\Omega$ is called \gesperrt{adjoining}, and one says that the field $\Omega'$ arises from $\Omega$ by \gesperrt{adjunction of} $\alpha$.

To $\Omega'$ one can again adjoin a magnitude $\alpha'$, and obtains a third field $\Omega''$, which then contains $\Omega$, $\alpha$ and $\alpha'$, and so on. But one can also adjoin two or more numbers to $\Omega$ simultaneously, and thus obtains the same field $\Omega''$ if one adjoins $\alpha$ and $\alpha'$ to $\Omega$ simultaneously. One can also adjoin to one field a second field, and thus obtains a field that contains the numbers of both fields at once. In the above example one would, by adjunction of the field $\Omega, \alpha'$ to the field $\Omega' = \Omega, \alpha$, again obtain the same field $\Omega''$.

If one adjoins to a field one of its parts, no new field arises.

Thus, for example, by adjoining to the field of rational numbers, which may be denoted by $\mathfrak{R}$, the number $i=\sqrt{-1}$
%CONT>>
